\documentclass[conference]{IEEEtran}
\usepackage[T1]{fontenc}
\usepackage[utf8]{inputenc}
\usepackage{times}
\usepackage{microtype}
\usepackage{enumitem}
\usepackage{booktabs}
\usepackage{hyperref}

\title{Control-Flow-First Hardware Construction:\\
Threads, Protocol APIs, and Backend Optimization}
\author{
\IEEEauthorblockN{Anonymous Submission}
\IEEEauthorblockA{OSCAR Workshop Extended Abstract}
}

\begin{document}
\maketitle

\begin{abstract}
Hardware design is still largely organized around structural commitment: modules, wires, and protocol state machines are chosen first, while control semantics emerge only indirectly from those structures. This abstract explores a different design center. We present a control-flow-first hardware construction style in which the primary frontend object is a \emph{thread}: a sequential control program over explicit actions such as \texttt{Step}, \texttt{waitCondition}, \texttt{jump}, \texttt{Call}, and \texttt{Return}. In this model, protocols are packaged as APIs rather than exposed as raw handshake choreography, and backend compilation is responsible for lowering control programs into concrete scheduling, pipelining, and resource-allocation decisions.

We argue three points through a Chisel-based prototype. First, hardware behavior can be written as control programs rather than structural interconnect skeletons, with \texttt{thread} serving as an explicit machine model over step-level timing units. Second, protocol-heavy interactions, such as AXI reads, can be encapsulated as reusable APIs because APIs are packaged control programs executed under the same thread semantics. Third, microarchitectural choices such as asynchronous path decoupling, write ordering, reservation discipline, and commit timing can be pushed into a backend/runtime layer without changing frontend API shape, as illustrated by a decoupled CPU case study and an age-ordered scoreboard regfile.
\end{abstract}

This extended abstract makes three contributions: (1) it defines \texttt{thread} as an explicit hardware execution model with first-class control-flow edges rather than a hidden runtime convenience; (2) it shows that protocol APIs can be understood as packaged control programs, using AXI as a concrete example; and (3) it shows that backend organization can move from synchronous linear execution chains to asynchronous call trees without forcing a corresponding change in frontend API shape.

\section{Control-Flow-First Hardware}
Our starting point is simple: a large fraction of hardware behavior is fundamentally sequential, yet it is commonly authored through structural interconnect first. We instead take \texttt{thread} as the primary frontend abstraction.

A \texttt{thread} is not a software thread, nor a hidden runtime mechanism. It is a sequential list of explicitly named timing units, \texttt{Step}s, together with a runtime cursor that advances among them. In that sense, a thread is a state machine with explicit default sequencing: absent a jump, return, or stall, control falls through from one step to the next. This makes \texttt{thread} an explicit machine model and \texttt{Step} a hardware control-flow instruction.

The resulting frontend is best viewed as a control-flow assembly for hardware. Unlike conventional software assembly, however, the model is not merely node-centric. It is a complete control-flow graph (CFG): \texttt{Step}s are programmable nodes, while return edges, jump edges, and patched edges are also first-class action sites.

\begin{table}[t]
\centering
\small
\begin{tabular}{@{}lll@{}}
\toprule
Dimension & Software assembly & This work \\
\midrule
Execution host & CPU core & \texttt{thread} \\
Program unit & instruction & \texttt{Step} \\
Control flow & node-centric stream & explicit CFG \\
Action attachment & node-local & node + edge \\
Default sequencing & PC+1 & next \texttt{Step} \\
Stall semantics & hidden & \texttt{waitCondition} \\
Reusable subroutine & function & API/program \\
\bottomrule
\end{tabular}
\caption{Software-assembly intuition versus control-flow-first hardware.}
\end{table}

\section{From Threads to Protocol APIs}
Once state-machine construction is generalized around \texttt{thread} and \texttt{Step}, the program being executed can be separated from the executor that runs it. A thread provides the execution model; a reusable control segment is then a packaged program meant to run on that model. In this sense, an API is not merely an interface wrapper around hardware. It is a packaged control program.

This interpretation is useful for protocols. Consider an AXI read. In a conventional style, each call site would directly manipulate \texttt{valid}/\texttt{ready} signals, response checking, and return sequencing. In our prototype, the AXI interaction is packaged as a callable control program. The user writes an API call; the implementation owns the internal handshake timing.

This works because API timing has two parts. Internal timing is \emph{closed}: the packaged program determines how the protocol proceeds internally. External timing is left \emph{composable}: the call site determines where the API is invoked, and the return site determines what happens when the API completes. In our framework, \texttt{ContextScope} resolves the call-site execution context, while continuations and edge actions resolve return behavior. A call may return to an explicit next \texttt{Step}, or it may trigger an edge-local guarded action. In this sense, an API call is itself a CFG-level object: it may be initiated from a node, but its completion semantics may be realized either as node-local continuation code or as edge-local actions.

\section{Backend-Driven Optimization}
The control-flow-first view becomes most useful when frontend code remains stable while backend realization changes. We illustrate this with a CPU-style case study and a regfile example.

In the CPU prototype, the important refactoring was not simply ``turning paths into processes.'' The key transformation was to make execution paths asynchronous while keeping the frontend API shape unchanged. Decode no longer owned the full downstream lifecycle of execution and writeback. Instead, decode became a router: it interpreted the instruction, selected a destination path, and launched an asynchronous transaction. Conceptually, this turns what would otherwise look like a linear function-call stack into an asynchronous call tree.

The same principle appears in the age-ordered scoreboard regfile. From the caller's perspective, frontend operations remain simple: reserve a destination, later write back and clear it. Internally, however, the backend maintains substantially richer policy: reservation tracking, pending-write buffering, ready-state publication, age ordering, and commit-time visibility. These are precisely the kinds of choices that should live below the frontend control description.

\begin{table}[t]
\centering
\small
\begin{tabular}{@{}ll@{}}
\toprule
Conventional notion & Control-flow-first reinterpretation \\
\midrule
Decode stage & decode router \\
Pipeline stage & async call segment \\
Backpressure & thread blocking \\
Stage handoff & async call into program \\
Writeback stage & commit service / continuation \\
Structural hazard & conflict at actual service boundary \\
\bottomrule
\end{tabular}
\caption{CPU concepts reinterpreted in a thread-centric model.}
\end{table}

\section{Takeaway}
This work does not argue that structure disappears. Hardware remains hardware, and concrete implementations still require wires, registers, and arbitration. Our claim is instead about factoring: control intent can be authored first, APIs can represent packaged protocol programs, and backend lowering can reintroduce structural machinery where it is actually needed.

The prototype is still evolving, and we do not yet claim a full optimizing compiler or final performance results. Nevertheless, the current artifact already suggests a useful direction for hardware construction: represent behavior as executable control programs, treat protocol interactions as APIs over those programs, and leave scheduling and microarchitectural policy to the backend.

\end{document}
